\begin{proposition}[Competition creates competition]\label{prop:ccc}
Let the classes satisfy the class ranking and have finite means, let \(F_1\le F_0\) and \(F_2\le F_1\) be
continuous incumbent laws, and fix \(0<\pi,\rho,\lambda<1\), \(0<\beta<1/2\), \(c_H>0\) and \(k>0\). Write
\(\Delta_i=\DT(F_i)\), \(B_i=B_{F_i}\), \(b=(1-\lambda)\beta\), and
\(\kappa(\pi)=\min\{m(\pi),1-M(\pi)\}\). Suppose that
\[
\begin{array}{ll}
\text{(H1), the high-cost window:}& B_0(\pi)<c_H<B_1(M(\pi)),\\
\text{(H2), the trading-cost window:}& \Delta_0<k<\rho\kappa(\pi)\Delta_1,\\
\text{(H3), the exclusion condition:}& B_2(M(\pi))<c_H.
\end{array}
\]
Each economy has a unique equilibrium outcome under all four tie rules.
\begin{enumerate}[label=(\roman*),leftmargin=2em]
\item At \(F_0\), both classes abstain, every posterior equals \(\pi\), entry is \(\rho\), and all orders share the
price \(t_{0,0}+\rho[t_{L,0}-t_{0,0}+\Delta_0\pi]\). High-class ownership is \(\pi\rho W_H(F_0)\).
\item At \(F_1\), class \(H\) buys and class \(L\) sells with probability one. The posteriors after
\((-1,0,1)\) are \((m(\pi),\pi,M(\pi))\), and entry is \((\rho,\rho,1)\). Prices rise strictly in this order.
Class entry and total entry are
\[
e_H=\rho+(1-\rho)(b+\lambda),\qquad e_L=\rho+(1-\rho)b,\qquad
\Ent_1=\rho+(1-\rho)(b+\pi\lambda)>\rho.
\]
High-class ownership rises exactly when \(e_HW_H(F_1)>\rho W_H(F_0)\). The probability that class \(H\) acquires
the target, conditional on acquisition by the challenger, rises exactly when
\begin{equation}\label{eq:a3-share}
e_HW_H(F_1)W_L(F_0)>e_LW_H(F_0)W_L(F_1).
\end{equation}
If class \(H\) outbids the stronger incumbent, \(W_H(F_1)=1\), it also outbids \(F_0\); high-class ownership
rises from \(\pi\rho\) to \(\pi e_H\), and its takeover share is nondecreasing, strictly exactly when
\(W_L(F_0)>0\). If \(W_L(F_0)=0\), that share is one at both laws.
\item At \(F_2\), the orders and posteriors are those of part (ii), entry is \(\rho\) after every order, and the
three prices differ. High-class ownership is \(\pi\rho W_H(F_2)\).
\end{enumerate}
At each law prices are \(P_i(x)=t_{0,i}+e_i(x)[t_{L,i}-t_{0,i}+\Delta_i\mu_i(x)]\), the price belief is
\(\mu_i(x)\), high-cost preparation is \([e_i(x)-\rho]/(1-\rho)\), and bids are truthful.
\end{proposition}
